\section{视觉的信仰：为什么选择 Computer Vision}

谢赛宁讲自己为什么选择视觉时，并没有从“这个方向更热门”或“更容易发论文”说起，而是从一个很孩子气、但极其根本的思想实验说起：如果必须去掉一个感官，哪个感官最不能失去？在他看来，视觉几乎无法替代。这种判断并不是简单的感官偏好，而是一种关于智能结构的直觉。一个系统如果无法从连续、模糊、带噪声的空间信号中提炼出可行动的世界表征，那么它距离真正的 intelligence 就还很远。也正因为如此，视觉并不是 AI 中的一个边缘应用，而更像理解智能的主入口之一。

他随后给出的论据带有典型的“从生物演化回看机器学习”的视角。视觉相关区域占大脑皮层的大约 \(30\%\)，如果把更广泛的激活范围算进去，甚至能到 \(70\%\)。这当然不是说“视觉就是全部”，而是说明视觉在人类智能架构里承担着超出单一感官的地位。谢赛宁引用寒武纪大爆发理论来解释这一点：当视觉出现之后，生物之间的竞争突然升级，因为“被看见”和“看见别人”同时发生了。这既提高了捕食与逃生的精度，也迫使整个进化系统转向更高阶的表征、预测和决策。

\begin{importantbox}{“眼睛是暴露在真实世界的大脑”}
访谈里最核心的一句判断是：\textit{``眼睛是唯一一个暴露在真实世界里面的大脑部分...解决视觉不是要解决视觉本身，而是要解决智能本身''}。这句话把视觉从一个“任务集合”提升成了一个“智能接口”。如果眼睛本质上是大脑伸向世界的前端，那么视觉研究的终点就不只是分类、检测或分割，而是学习如何把世界压缩成可用于理解、预测和行动的内部表征。
\end{importantbox}

因此，谢赛宁对 Computer Vision 的定义天然带有反潮流的色彩。他不愿把 vision 视为某个 benchmark 的集合，也不愿把它仅仅理解成图像任务、自动驾驶或多模态系统的某个组件。在他看来，Computer Vision 更接近一种 perspective，一种处理真实世界信息的方式。这个 perspective 至少包含两层含义：第一，视觉面对的是连续空间、高维度且高度带噪声的输入，不能像语言那样天然享受离散 token 化带来的结构红利；第二，视觉必须依赖层次化表征，从局部纹理、边界、形状一路上升到物体、场景、关系与动态因果。也正因为这两层特征，视觉从来都不只是“看图说话”，而是智能系统如何与物理现实接壤的集中体现。

\begin{knowledgebox}{从寒武纪到深度视觉}
谢赛宁之所以反复讲寒武纪大爆发，不是为了做生物学类比，而是为了说明一件事：视觉一旦出现，整个智能系统的复杂度就会被重新定义。对于机器学习来说，视觉问题之所以难，不是因为像素多，而是因为像素背后对应的是三维空间、对象恒常性、遮挡、运动、因果和行动后果。这些结构远比 token 预测要厚重。
\end{knowledgebox}

这种“视觉是北极星”的看法，也解释了他为什么在大模型浪潮里并不沮丧。很多做视觉的人会担心，LLM 的进展会不会让视觉沦为语言模型的附庸；但在谢赛宁那里，答案几乎是相反的。正因为语言模型证明了表征学习与规模化训练的威力，视觉研究者才更有机会重新争论“真正的瓶颈是什么”。语言的强势，并不意味着世界被语言穷尽了；恰恰相反，它暴露出一个更大的缺口：大量真实世界知识并不能被自然语言完整编码，更不能靠语言在数字空间中的自回归过程自动补上。

他对语言与视觉关系的态度因此非常平衡。语言很重要，因为它给问题定义、抽象和人机对齐提供了强大工具；但如果让语言成为唯一接口，就容易把系统困在一个“说得明白但摸不到世界”的层面。视觉在这里不是为了给语言打工，而是为了重新打开与真实世界的连接。后来他把 LLM 称为“virtual intelligence”，也正是基于这层区分：语言模型能在数字世界里完成很多惊人的操作，但要面对空间、身体、行动和物理约束，它还需要一种更深的世界表征作为底座。

\begin{warningbox}{把 Vision 缩减成 benchmark，会错过它真正的野心}
如果把 Computer Vision 理解成“目标检测、图像分类、视频理解”的若干竞赛赛道，就会自然得出一个错误结论：这些任务迟早会被更大的 foundation model 吸收，所以视觉会失去独立价值。谢赛宁的观点恰恰是，视觉之所以重要，不在于这些任务本身，而在于它们共同指向同一个核心难题：怎样从复杂、连续、具身的世界里学习出稳定而可泛化的表征。
\end{warningbox}

值得注意的是，谢赛宁对视觉的执着，并不来自学科忠诚，而更像是一种“认知假设”：如果一个系统无法处理大规模连续信号、无法在多层表征之间形成稳定映射、无法把感知结果用于后续预测与规划，那么它就很难被称为真正意义上的智能。换言之，视觉吸引他的不是图像本身，而是图像背后那整套关于空间、遮挡、时间和行动可供性的复杂性。也因此，做视觉在他那里从来不是进入一个较窄的子领域，而是主动选择去做最难被离散化、也最接近现实世界的一类问题。

从这个角度回头看，他后来对长镜头、世界模型和 wearable/robotics 的兴趣，其实早已被写在这一最初选择里。一个少年凭直觉觉得“最不能失去视觉”，后来成为一个研究者继续追问“为什么视觉在智能里如此中心”，中间看似跨越多年，内核却始终没有改变。这种稳定的问题意识，恰恰比具体技术名词更值得记录。

从这个意义上说，谢赛宁选择视觉，并不是“选择一个方向”，而是“选择一种关于智能的信仰”。这种信仰让他后来无论做自监督、生成模型、多模态，还是世界模型，始终没有真的离开过 Computer Vision。因为在他的语境里，vision 从来不是一个即将被替代的子领域，而是智能研究不可绕开的主战场。

\subsection{本章小结}

谢赛宁之所以坚定地选择 Computer Vision，是因为他把视觉视为智能与真实世界之间的核心接口，而不是若干视觉任务的总和。寒武纪视角、脑科学比例、层次化表征和“眼睛是暴露在真实世界的大脑”这几个判断，共同构成了他后续所有研究路线的理论底座。

